%% --- PACKAGES PROPRES AUX COURS ---
\usepackage{fancyhdr}
\usepackage{lastpage}
\usepackage[column=X]{cellspace}

%% --- VARIABLES ---
\newcommand{\monTitre}{Titre du cours} % À redéfinir dans chaque fichier : \renewcommand{\monTitre}{Mon titre}

%% --- FORMAT DES TITRES (titlesec) ---

\titleformat{\chapter}[frame]
{\color{red!85!white}\normalsize}
{\filright\sffamily\Large\enspace Chapitre \thechapter\enspace}
{8pt}
{\color{red!85!white}\rule{0pt}{30pt}\sffamily\Huge\bfseries\filcenter}
\titlespacing{\chapter}{0pc}{0pc}{2pc}

\titleformat{\section}
{\color{green!55!black}\normalfont\Large\bfseries}
{\thesection}
{1em}
{}

% 3. Subsections (Cyan)
% J'ai retiré le \hspace{1em} qui poussait le texte
\titleformat{\subsection}[block]
{\color{cyan!65!black}\normalfont\large\bfseries}
{\thesubsection}
{1em}
{}
% J'ai mis le premier paramètre à 0pt pour supprimer l'indentation de bloc
\titlespacing{\subsection}{0pt}{1.5ex}{1ex}

% 4. Subsubsections (Violet)
\titleformat{\subsubsection}[block]
{\color{violet}\normalfont\normalsize\bfseries}
{\thesubsubsection}
{1em}
{}
\titlespacing{\subsubsection}{0pt}{1ex}{1ex}

%% --- NUMÉROTATION ---
\renewcommand{\thesection}{\Roman{section}.}
\renewcommand{\thesubsection}{\arabic{subsection})}
\renewcommand{\thesubsubsection}{\alph{subsubsection})}
\setcounter{secnumdepth}{3}

%% --- EN-TÊTES ET PIEDS DE PAGE ---
\pagestyle{fancy}
\renewcommand\headrulewidth{1pt}
\fancyhead[LO,RE]{\maClasse}
\fancyhead[RO,LE]{\monTitre}
\renewcommand\footrulewidth{1pt}
\fancyfoot[LO,RE]{Année 2025-2026}
\fancyfoot[RO,LE]{\textbf{Page \thepage/\pageref{LastPage}}}

%% --- RÉGLAGES PARAGRAPHES ---
\setlength{\cellspacebottomlimit}{4pt}
\setlength{\cellspacetoplimit}{4pt}
\setlength{\parindent}{2pt}
\setlength{\tabcolsep}{0.5cm}

%% --- COMPTEURS ET BOITES TCOLORBOX ---

%% --- COMPTEURS AVEC RÉINITIALISATION À CHAQUE CHAPITRE ---

% On ajoute [chapter] à la fin pour que le compteur reparte à 0 à chaque \chapter
\newcounter{dfn}[chapter]
\newcounter{thm}[chapter]
\newcounter{cvn}[chapter]
\newcounter{prp}[chapter]
\newcounter{ex}[chapter]
\newcounter{mtd}[chapter]
\newcounter{rq}[chapter]
\newcounter{demo}[chapter]
\newcounter{appli}[chapter]
\newcounter{exo}[chapter]
\newcounter{lem}[chapter]

%% --- BLOC COMPLET DES BOITES TCOLORBOX ---

% --- DÉFINITIONS (Bleu) ---
\newtcolorbox{definition}{breakable,colback=white!6!white,colframe=blue!85!white,fonttitle=\bfseries,title=Définition \thedfn, code={\stepcounter{dfn}}, arc=3mm}
\newtcolorbox{definitionTitre}[1]{breakable,colback=white!6!white,colframe=blue!85!white,fonttitle=\bfseries,title=Définition \thedfn\ - #1, code={\stepcounter{dfn}}, arc=3mm}
\newtcolorbox{definitions}{breakable,colback=white!6!white,colframe=blue!85!white,fonttitle=\bfseries,title=Définitions \thedfn, code={\stepcounter{dfn}}, arc=3mm}

% --- PROPRIÉTÉS & THÉORÈMES (Rouge/Orange) ---
\newtcolorbox{theoreme}{breakable,colback=red!6!white,colframe=red!85!white,fonttitle=\bfseries,title=Théorème \thethm, code={\stepcounter{thm}}, arc=3mm}
\newtcolorbox{theoremeTitre}[1]{breakable,colback=red!6!white,colframe=red!85!white,fonttitle=\bfseries,title=Théorème \thethm\ - #1, code={\stepcounter{thm}}, arc=3mm}
\newtcolorbox{propriete}{breakable,colback=red!6!white,colframe=red!85!white,fonttitle=\bfseries,title=Propriété \theprp, code={\stepcounter{prp}}, arc=3mm}
\newtcolorbox{proprieteTitre}[1]{breakable,colback=red!6!white,colframe=red!85!white,fonttitle=\bfseries,title=Propriété \theprp\ - #1, code={\stepcounter{prp}}, arc=3mm}
\newtcolorbox{proprietes}{breakable,colback=red!6!white,colframe=red!85!white,fonttitle=\bfseries,title=Propriétés \theprp, code={\stepcounter{prp}}, arc=3mm}
\newtcolorbox{lemme}{breakable,colback=red!6!white,colframe=orange!85!white,fonttitle=\bfseries,title=Lemme \theprp, code={\stepcounter{prp}}, arc=3mm}
\newtcolorbox{convention}{breakable,colback=orange!6!white,colframe=orange!85!white,fonttitle=\bfseries,title=Convention \thecvn, code={\stepcounter{cvn}}, arc=3mm}

% --- EXEMPLES (Vert) — argument optionnel pour le titre ---
% Usage : \begin{exemple}             -> "Exemple 1"
%         \begin{exemple}[Mon titre]  -> "Exemple 1 - Mon titre"
\newtcolorbox{exemple}[1][]{
  enhanced,
  breakable,
  nobeforeafter,
  % --- ON FORCE TOUT À ZÉRO ---
  boxrule=0pt,        % Épaisseur du cadre global à 0
  titlerule=0pt,      % Supprime la ligne sous le titre
  sharp corners,      % Important pour les coupures propres
  % ----------------------------
  colback=white,
  colframe=white,     % On met le cadre en blanc (invisible) au cas où
  frame empty,        % On cache le cadre
  %
  borderline west={0.5mm}{0mm}{green!55!black}, % Seule la barre verte reste
  %
  % Gestion de la coupure
  toprule at break=0pt,
  bottomrule at break=0pt,
  %
  title={Exemple \theex\ifx&#1&\else\ - #1\fi},
  code={\stepcounter{ex}},
  attach boxed title to top left = {xshift = -3mm},
  boxed title style={size=small,colback=green!55!black,colframe=green!55!black, arc=3mm}
}

% --- REMARQUES (Orange) ---
\newtcolorbox{remarque}{enhanced, breakable, frame empty, nobeforeafter, borderline west={0.5mm}{0mm}{orange!85!white}, title=Remarque \therq, colback=white, code={\stepcounter{rq}}, attach boxed title to top left = {xshift = -3mm}, boxed title style={size=small,colback=orange!85!white,colframe=orange!85!white, arc=3mm}}
\newtcolorbox{remarques}{enhanced, breakable, frame empty, nobeforeafter, borderline west={0.5mm}{0mm}{orange!85!white}, title=Remarques \therq, colback=white, code={\stepcounter{rq}}, attach boxed title to top left = {xshift = -3mm}, boxed title style={size=small,colback=orange!85!white,colframe=orange!85!white, arc=3mm}}
\newtcolorbox{remarqueTitre}[1]{enhanced, breakable, frame empty, nobeforeafter, borderline west={0.5mm}{0mm}{orange!85!white}, title=Remarque \therq\ - #1, colback=white, code={\stepcounter{rq}}, attach boxed title to top left = {xshift = -3mm}, boxed title style={size=small,colback=orange!85!white,colframe=orange!85!white, arc=3mm}}

% --- MÉTHODES & APPLICATIONS (Bleu foncé) ---
\newtcolorbox{methode}{enhanced, breakable, frame empty, nobeforeafter, borderline west={0.5mm}{0mm}{blue!75!black}, title=Méthode \themtd, colback=white, code={\stepcounter{mtd}}, attach boxed title to top left = {xshift = -3mm}, boxed title style={size=small,colback=blue!75!black,colframe=blue!75!black, arc=3mm}}
\newtcolorbox{methodes}{enhanced, breakable, frame empty, nobeforeafter, borderline west={0.5mm}{0mm}{blue!75!black}, title=Méthodes \themtd, colback=white, code={\stepcounter{mtd}}, attach boxed title to top left = {xshift = -3mm}, boxed title style={size=small,colback=blue!75!black,colframe=blue!75!black, arc=3mm}}
\newtcolorbox{methodeTitre}[1]{enhanced, breakable, frame empty, nobeforeafter, borderline west={0.5mm}{0mm}{blue!75!black}, title=Méthode \themtd\ - #1, colback=white, code={\stepcounter{mtd}}, attach boxed title to top left = {xshift = -3mm}, boxed title style={size=small,colback=blue!75!black,colframe=blue!75!black, arc=3mm}}
\newtcolorbox{application}{enhanced, breakable, frame empty, nobeforeafter, borderline west={0.5mm}{0mm}{orange!75!black}, title=Application \themtd, colback=white, code={\stepcounter{mtd}}, attach boxed title to top left = {xshift = -3mm}, boxed title style={size=small,colback=orange!75!black,colframe=orange!75!black, arc=3mm}}

% --- DIVERS (Exercices, Python, Conséquences) ---
\newtcolorbox{exercice}{enhanced, breakable, frame empty, nobeforeafter, borderline west={0.5mm}{0mm}{green!55!black}, title=Exercice n°\theexo, colback=white, code={\stepcounter{exo}}, attach boxed title to top left = {xshift = -3mm}, boxed title style={size=small,colback=green!55!black,colframe=green!55!black, arc=3mm}}
\newtcolorbox{exerciceTitre}[1]{enhanced, breakable, frame empty, borderline west={0.5mm}{0mm}{green!55!black}, title=Exercice n°\theexo\ - #1, colback=white, code={\stepcounter{exo}}, attach boxed title to top left = {xshift = -3mm}, boxed title style={size=small,colback=green!55!black,colframe=green!55!black, arc=3mm}}
\newtcolorbox{consequences}{enhanced, breakable, frame empty, nobeforeafter, borderline west={0.5mm}{0mm}{ProcessBlue}, title=Conséquences, colback=white, attach boxed title to top left = {xshift = -3mm}, boxed title style={size=small,colback=ProcessBlue,colframe=ProcessBlue, arc=3mm}}
\newtcolorbox{programme}{enhanced, breakable, frame empty, nobeforeafter, borderline west={0.5mm}{0mm}{gray!55!black}, title=Programme Python, colback=white, attach boxed title to top left = {xshift = -3mm}, boxed title style={size=small,colback=gray!55!black,colframe=gray!55!black, arc=3mm}}


% Commande pour les démonstrations (incrémentation automatique)
\newcommand{\demonstration}{%
  \stepcounter{demo}%
  \tcbox[enhanced, 
         on line, 
         arc=2mm, 
         colback=gray!10!white, 
         colframe=gray!50!black, 
         size=small, 
         fontupper=\bfseries]
         {Démonstration \thedemo}%
  \enspace % Ajoute un petit espace après la boîte
}

%% --- FIN DU BLOC DES BOITES ---

% Activité (Numéro, Titre, Chapitre)
\newcommand{\enteteActivite}[3]{
    \noindent
    \begin{tabularx}{\textwidth}{@{}p{2cm} >{\centering\arraybackslash}X p{2cm}@{}}
        \textbf{Activité n°#1} & \textbf{#2} & \textbf{Chapitre #3} \\
    \end{tabularx}
}


% --- Commande pour les blocs mémo personnalisés ---
% Argument 1 : Couleur du cadre
% Argument 2 : Titre du bloc
% Argument 3 : Texte/Contenu
\newcommand{\blocMemo}[3]{%
    \begin{tcolorbox}[
        enhanced,
        breakable,
        title=#2,
        colback=#1!5,          % Fond très clair
        colframe=#1!75!black,  % Bordure
        fonttitle=\bfseries\sffamily,
        colbacktitle=#1!75!black,
        attach boxed title to top left={xshift=3mm, yshift=-3mm},
        boxed title style={size=small, arc=2mm, colframe=#1!75!black},
        arc=3mm,               % Bel arrondi
        boxrule=0.6mm,         % Trait un peu plus épais
        after skip=15pt,
        before skip=10pt
    ]
        \vspace{2mm} % Petit espace pour ne pas coller au titre boxé
        #3
    \end{tcolorbox}
}


% --- Nouvel environnement Culture Maths ---
\newtcolorbox{cultureMaths}{
    enhanced, 
    breakable, 
    frame empty, 
    nobeforeafter, 
    borderline west={0.5mm}{0mm}{gray!60!black}, % Bordure grise comme ton "programme"
    title=Le saviez-vous ?, 
    colback=white, 
    attach boxed title to top left = {xshift = -3mm}, 
    boxed title style={size=small,colback=gray!60!black,colframe=gray!60!black, arc=3mm},
    coltitle=white,
    fonttitle=\bfseries\sffamily,
    after skip=15pt,
    before skip=10pt
}